% Generated draft: requires mathematical and source review.
\documentclass[11pt]{article}
\usepackage[margin=25mm]{geometry}
\usepackage{fontspec}
\setmainfont{Latin Modern Roman}
\usepackage{amsmath,amssymb,mathtools,unicode-math}
\setmathfont{Latin Modern Math}
\usepackage{enumitem,needspace,xurl,hyperref}
\setlength{\parindent}{0pt}
\setlength{\parskip}{5pt}
\setlength{\emergencystretch}{3em}
\newcommand{\N}{\mathbb{N}}
\newcommand{\Z}{\mathbb{Z}}
\newcommand{\Q}{\mathbb{Q}}
\newcommand{\R}{\mathbb{R}}
\newcommand{\C}{\mathbb{C}}
\newcommand{\F}{\mathbb{F}}
\newcommand{\A}{\mathbb{A}}
\newcommand{\PP}{\mathbb{P}}
\newcommand{\E}{\mathbb{E}}
\newcommand{\eps}{\varepsilon}
\newcommand{\dd}{\,\mathrm{d}}
\begin{document}
\section*{UnsolvedMath ID: 30001919}
\textbf{Source problem code:} catalogue entry
\subsection{Definitions and mathematical statement}
Let \(X\) be a Fano manifold, so that its anticanonical line bundle \(-K_X\) is ample, and let \(D\in |-K_X|\) be a smooth divisor. Write \(c_1(X)=c_1(-K_X)\). Define the greatest Ricci lower bound by
\[
R(X)=\sup\bigl\{t\in\R:\text{ there exists a smooth Kähler form }\omega\in c_1(X)\text{ with }\operatorname{Ric}(\omega)-t\omega>0\bigr\},
\]
where the inequality means positivity as a real \((1,1)\)-form. Use compatible normalizations in which \(\operatorname{Ric}(\omega)\) represents \(c_1(X)\), and let \(\{D\}\) denote the current of integration along \(D\), normalized to represent \(c_1(\mathcal O_X(D))=c_1(X)\).

For \(0<\beta<1\), the usual interpretation of a cone-singularity solution is a Kähler metric \(\omega_\beta\) on \(X\setminus D\), extending to a closed positive \((1,1)\)-current in \(c_1(X)\), with cone angle \(2\pi\beta\) along \(D\), satisfying
\[
\operatorname{Ric}(\omega_\beta)
=\beta\omega_\beta+(1-\beta)\{D\}
\]
as an equality of currents on \(X\). Locally, in holomorphic coordinates with \(D=\{z_1=0\}\), the cone-angle condition is ordinarily expressed by uniform comparability on the complement of \(D\) with
\[
\sqrt{-1}\left(|z_1|^{2\beta-2}\,\dd z_1\wedge\dd\overline z_1+
\sum_{j=2}^{\dim_{\C}X}\dd z_j\wedge\dd\overline z_j\right).
\]
The supplied record does not specify the precise regularity class, the definition of the singular Ricci current, or normalization factors. The preceding conventions explain the standard intended interpretation; the exact admissible solution class requires confirmation from the source.

\textbf{Source conjecture.} For every such pair \((X,D)\), does a cone-singularity solution exist for every real \(\beta\) satisfying \(0<\beta<R(X)\)? Moreover, if \(R(X)<1\), is there no such solution for any real \(\beta\) satisfying \(R(X)<\beta<1\)? The supplied formulation makes no assertion at \(\beta=R(X)\), \(\beta=0\), or \(\beta=1\). These are conjectural assertions as supplied, rather than results established by this draft.
\subsection{Short English statement}
For a Fano manifold \(X\) and a fixed smooth anticanonical divisor \(D\), the conjecture asks whether conical Kähler--Einstein metrics exist at every angle \(2\pi\beta\) below the threshold \(R(X)\), and fail to exist above that threshold when \(R(X)<1\). The threshold itself is excluded. The precise analytic meaning of a cone-singularity solution needs source confirmation.
\subsection{Sources}
\begin{itemize}
\item Supplied local UnsolvedMath registry record, permanent ID \(30001919\), source problem code \texttt{catalogue entry}, titled ``Conical Kähler--Einstein Metrics and Ricci Lower Bounds,'' in the Oberwolfach Reports Open Problems collection. No individual proposer is supplied; the proposed year is recorded as 2011.
\item Supplied report citation: \emph{Complex Algebraic Geometry (2012)}, Oberwolfach Reports, DOI 10.4171/owr/2011/47, pp.\ 2669--2733. \url{https://doi.org/10.4171/owr/2011/47}. The background describes these reports as documenting mathematical workshops and their open-problem discussions. The citation's 2012 designation and the record's 2011 year information are retained without reconciliation.
\item Supplied literature reference: Song--Wang, ``greatest Ricci lower bound and conical metrics.'' \url{https://arxiv.org/abs/1207.4839}.
\item Supplied literature reference: Chen--Donaldson--Sun, ``conical Kähler--Einstein metric revisited.'' \url{https://arxiv.org/abs/1207.5011}.
\item Literature assessment dated exactly \texttt{2026-08-21}: the supplied curation labels the problem ``partially solved'' and states that the conjecture is proved in toric and other substantial settings, including existence below the greatest Ricci lower bound under additional hypotheses, but not in the full generality stated. This is supplied dated curation, not independently verified here. The record expressly calls for independent verification of time-sensitive status and cited claims.
\item The record labels its statement \texttt{corrected\_verified} and reports that the mathematical content was checked against the official Oberwolfach report, with notation, context, and question form normalized without changing the source-level claim. This verification is reported by the supplied record and has not been independently performed for this draft.
\end{itemize}
\end{document}
